\documentclass[11pt]{article}
\usepackage{fontspec}
\usepackage[french]{babel}
\usepackage{amsmath}
\usepackage{amssymb}
\usepackage{amsthm}
\usepackage{mathrsfs}
\usepackage{tikz-cd}
\usepackage{enumitem}

\newcommand{\printedpage}[1]{\marginpar{\footnotesize p.~#1}}
\newcommand{\msp}[1]{\marginpar{\footnotesize [#1]}}
\newcommand{\circled}[1]{\tikz[baseline=(c.base)]\node[draw,circle,inner sep=1pt](c){#1};}
\newcommand{\Ob}{\operatorname{Ob}}

\begin{document}

\printedpage{58}
% head: A. Grothendieck : Dérivateurs, Ch. XII, Édition des Universités de Montpellier II et Paris VII
% foot: Version 7 décembre 2002

est \textit{surjectif.} (P. ex. si $S'=\emptyset$, on a $X'=\emptyset$, donc $f'\in W$, sans que ça implique $f\in W$!) Le fait que dans le cas cocartésien de W(7 bis), $i_0\in W\Longrightarrow i\in W$ doit être considéré comme ‘allant de soi’ - autrement le carré ne mérite pas, du point de vue $W$-homotopique, le nom de ‘cocartésien’. Par contre, l’implication inverse $i\in W\Longrightarrow i_0\in W$ m’apparaît comme un trait très particulier du dérivatuer HOT$_W$ (quand il le possède), et notamment % [sic] dérivatuer -> dérivateur, misprint
de HOT $=$ HOT$_{W_\infty}$. Le début
\msp{page 80}
de la remarque signifie simplement que le dérivateur dual HOT$^\circ$ ne possède pas cette même propriété. Je dois renvoyer à ce sujet aux deux axiomes pour les dérivateurs correspondants à W(7 bis), W(7), ce que je ferai sans doute sous peu, de façon à donner un sens plus précis aux présents commentaires.

\subsection*{2.2 Application aux carrés cocartésiens déduits d’une immersion ouverte ou fermée.}

Les axiomes W(7), W(7 bis) (du cube, et du carré cocartésien) concernent des carrés cocartésiens de nature très particulière. Nous allons voir à présent comment en déduire des énoncés du même type pour des carrés cocartésiens beaucoup plus généraux, y compris pour des carrés qui n’ont rien de cocartésien (dans Cat) - et que nous appelerons % [sic] appelerons -> appellerons, misspelling
‘homotopiquement cocartésiens’ (relativement à $W$), ou simplement $W$-cocartésiens. (Il y aura également des carrés $W$-cartésiens un jour \dots) Dans le cas où $W=W_\infty$, je les appelerai aussi Hot-cocartésiens. % [sic] appelerai -> appellerai, misspelling

Je commencerai par le cas particulier d’un carré cocartésien associé à une immersion ouverte (diagramme ($*$), p. 74) :
\[
(*)\qquad\qquad
\begin{tikzcd}
U \arrow[r, hook, "i"] \arrow[d, "p"'] & X \arrow[d, "q"'] & T \arrow[l, hook', dotted, "j"'] \arrow[d, dotted] \\
U' \arrow[r, hook, "i'"] & X' & T'\simeq T. \arrow[l, hook', dotted, "j'"']
\end{tikzcd}
\]
\msp{page 81}

Le grand traquenard ici qu’il faut arriver à exorciser, c’est qu’il peut arriver que $i\in W$, sans qu’on ait $i'\in W$, et aussi $p\in W$ sans qu’on ait $q\in W$. Je m’en étais convaincu en 1983 en écrivant Pursuing Stacks (durant les semaines de ‘scratchwork’), et n’ai pas envie de rechercher un contre-exemple ici, s’il ne se présente de lui-même (ce qu’il fera sans doute).

Je vais donner un cas important où ces ennuis disparaissent :

\textbf{Proposition 6 :} \textit{Considérons les adhérences $\bar U$ et $\overline{U'}$ de $U$ et $U'$ dans $X$, $X'$ respectivement, et supposons que les inclusions}
\[
U\xrightarrow{i_0}\bar U,\qquad U'\xrightarrow{i'_0}\overline{U'}
\]
\printedpage{59}
% head: A. Grothendieck : Dérivateurs, Ch. XII, Édition des Universités de Montpellier II et Paris VII
% foot: Version 7 décembre 2002
\textit{soient dans $W$} (\footnote[64]{cf. th. 3 p. 107 pour des conditions générales assurant que cette hypothèse est satisfaite.}). \textit{Alors}
\[
\left\{
\begin{array}{llll}
\circled{\textit{a}}\ i\in W & \iff & i'\in W & \textit{si on a }\mathrm{W}(7\text{ bis})\\
\circled{\textit{b}}\ p\in W & \iff & q\in W & \textit{si on a }\mathrm{W}(7).
\end{array}
\right.
\]

\textsc{Démonstration}. Considérons d’abord le cas \circled{a} de $i$, $i'$. Soit
\[
k:\bar U\hookrightarrow X
\]
l’inclusion. Le triangle
\begin{tikzcd}
U \arrow[r, "i_0"] \arrow[rr, bend right=8, "i"'] & \bar U \arrow[r, "k"] & X
\end{tikzcd}
, où $i_0\in W$, montre que
\[
i\in W\iff k\in W.
\]

D’autre part, $X$ est la réunion des fermés $\bar U$, $T$, donc on a un carré cocartésien dans Cat
\[
(*)\qquad\qquad
\begin{tikzcd}
\bar U\cap T=\dot T \arrow[r, "k_0"] \arrow[d, "j_0"'] & T \arrow[d, "j"] \\
\bar U \arrow[r, "k"] & X.
\end{tikzcd}
\]

L’axiome W(7 bis) (sous sa forme duale (i.e. appliqué au carré formé par les catégories opposées \dots)) implique
\msp{page 82}
que
\[
k\in W\iff k_0\in W,
\]
d’où
\[
i\in W\iff (k_0:\dot T\longrightarrow T)\in W.
\]
De même on trouve, en appliqunt ceci à $(X',U')$ au lieu de $(X,U)$ % [sic] appliqunt -> appliquant, misprint
\[
i'\in W\iff (k'_0:\dot T'\longrightarrow T')\in W.
\]
Or on a $T\xrightarrow{\sim}T'$ (isomorphisme), et par la proposition 4 cet isomorphisme transforme $\dot T$ en $\dot T'$. Donc
\[
k_0\in W\iff k'_0\in W\qquad\quad\text{(trivialement)},
\]
d’où a).

Prouvons b). Le carré commutatif induit par ($*$), p. 80,
\printedpage{60}
% head: A. Grothendieck : Dérivateurs, Ch. XII, Édition des Universités de Montpellier II et Paris VII
% foot: Version 7 décembre 2002
\[
\begin{tikzcd}
U \arrow[r, hook, "i_0"] \arrow[d, "p"'] & \bar U \arrow[d, "\bar p"] \\
U' \arrow[r, hook, "i'_0"] & \overline{U'},
\end{tikzcd}
\]
où $i_0,i'_0\in W$ par hypothèse, montre que
\[
p\in W\iff\bar p\in W,
\]
donc tout revient à prouver que
\[
\bar p\in W\overset{?}{\iff}q\in W.
\]
Mais $q$ induit un homomorphisme du carré ($*$) déduit de $(X,\bar U,T)$ (avec $X=\bar U\cup T$), dans le carré similaire pour $X'$
\[
\begin{tikzcd}
\overline{U'}\cap T'=\dot T' \arrow[r, "k'_0"] \arrow[d, "j'_0"'] & T \arrow[d, "j'"] \\ % [sic] T -> T', the square for X' should have the corner T' (although T'\simeq T)
\overline{U'} \arrow[r, "k'"] & X',
\end{tikzcd}
\]
les flèches verticales étant notées $q_{\dot T}$, $q_T$, $\bar p$, $q$.
\msp{page 83}
On sait que $q_{\dot T}$, $q_T$ sont des isomorphismes, donc dans $W$, donc le ‘lemme des cinq’ W(7) nous assure que $\bar p\in W\iff q\in W$, q.e.d.

\textbf{Commentaire.} Cette proposition nous amène à chercher des conditions sur un ouvert $U$ d’un $X$, qui assurent que l’inclusion $i_0:U\longrightarrow\bar U$ soit \texttt{[dans]} $W$, et \textit{que cette propriété soit conservée par tout cochangement de base $U\longrightarrow U'$.} De plus, nous recherchons des conditions qui soient de nature locale sur $\bar U$, plus exactement sur $\bar U\cap T=\dot U=\dot T$. Nous pouvons motiver ce souci (a posteriori) en disant que nous voudrions que la propriété $i_0\in W$ soit conservée, non seulement par cochangement de base quelconque $U\longrightarrow U'$, mais aussi, de plus, par passage de $X$ à toute sous-catégorie pleine $X_1$ contenant $U$, donnant donc naissance à une situation
\[
U\overset{i_Y}{\hookrightarrow}X_1,\qquad U\xrightarrow{(i_Y)_0}\bar U^{X_1}=\bar U\cap X_1
\]
\texttt{[}$Y=X_1$\texttt{?]}, similaire à la situation de départ. Plus généralement encore, nous voudrions qu’elle soit conservée en passant du triple $(X,U,T)$ à un triple “induit” $(X_1,U,T_1)$, par une flèche quelconque $g:T_1\longrightarrow T$ de but $T$ dans Cat, donnant naissance au couple de flèches
\[
U\xrightarrow{\mathrm{id}}U,\ g:T_1\longrightarrow T
\]
(cf. pages 60, 61). Le cas précédent est celui où
\printedpage{61}
% head: A. Grothendieck : Dérivateurs, Ch. XII, Édition des Universités de Montpellier II et Paris VII
% foot: Version 7 décembre 2002
\msp{page 84}
$T_1$ est une sous-catégorie pleine de $T$, $g$ la flèche d’inclusion, Notons à ce sujet la proposition soritale % [sic] inclusion, Notons -> inclusion. Notons, comma before a capital letter

\textbf{Proposition 7.} \textit{Conditions équivalentes sur l’ouvert $U$ d’un objet $X$ de} Cat (moyennant W(7 bis)) :

\begin{enumerate}[label={}, leftmargin=*]
\item[\textit{a)}] \textit{L’inclusion $i_0:U\longrightarrow\bar U$ est $W$-asphérique, i.e. pour tout $t\in\dot U=\bar U\backslash U$, la catégorie $U/t$ est $W$-asphérique.}
\item[\textit{b)}] \textit{L’inclusion $i_0:U\longrightarrow\bar U$ est dans $W$, et cette condition reste satisfaite, quand on remplace $(X,U)$ par $(X_1,U)$, où $(X_1,U)$ est déduit de $(X,U)$ au moyen d’une flèche $T_1\longrightarrow T(=X\backslash U)$ dans} Cat\textit{, comme rappelé ci-dessus.}
\item[\textit{b$'$)}] \textit{Comme b), dans le cas où $T_1$ est la catégorie ponctuelle $e$, et où de plus son image est dans $\dot U$.}
\end{enumerate}

\textsc{Démonstration}. Il est évident que si dans b), b$'$) on remplace la condition “$i_0\in W$” par la condition plus forte ‘$i_0$ est $W$-asphérique’, les quatre \texttt{[plutôt trois]} conditions a) à b$'$) sont équivalentes. Il reste à prouver que la condition que les $U/t$ soient $W$-asphériques ($t\in\bar U$ - mais c’est trivial pour $t\in U$, donc il suffit de l’exprimer pour $t\in\bar U\backslash U=\dot U$) est déjà impliquée par la plus faible b$'$ et on devine qu’on a ceci :

\msp{page 85}

\textbf{Corollaire.} (Moyennant W(7 bis).) \textit{Soit $X$ dans} Cat\textit{, $U$ un ouvert de $X$ tel que $X\backslash U=\overline{\{t\}}$ soit la catégorie finale. Alors l’inclusion $U\xrightarrow{i}X$ est dans $W$ si et seulement si elle est $W$-asphérique, i.e. si et seulement si $U/t$ est $W$-asphérique (où $t$ est l’unique élément de $T$).}

Nous démontrons ultérieurement ce corollaire (\footnote[65]{cf. cor. 4 de la proposition 15, page 149.}). La proposition 7 ne nous servira pas d’ici là - elle n’aura qu’un rôle purement esthétique à jouer, pour justifier notre propos délibéré de chercher des conditions sur $i:U\longrightarrow X$ (pour assurer que $U\longrightarrow\bar U$ est dans $W$ et le reste, quand on remplace $(X,U)$ par d’autres couples $(X_1,U_1)$ qui s’en déduisent par certaines opérations \dots) qui
\msp{page 86}
soient de \textit{nature locale} sur $\bar U$ (en fait, en les points de $\dot U$), et qui en tout état de cause entrainent la \textit{$W$-asphéricité} de $U\longrightarrow\bar U$.

J’avais espéré que cette condition d’asphéricité (suggérée par la notion ‘d’immersion de Dwyer’) serait stable par cochangement de base quelconque $U\longrightarrow U'$, ce qui aurait assuré en même temps qu’après une telle opération, l’immersion ouverte $i':U'\longrightarrow X'$ obtenue satisfait encore à la condition que $i'_0:U'\longrightarrow\overline{U'}$ est dans $W$. Nous verrons que cette attente n’est pas fondée - la condition de $W$-asphéricité n’est \textit{pas} stable par cochangement de base, et elle ne suffit \textit{pas} même à assurer qu’après un cochangement de base, $i'_0:U'\longrightarrow\overline{U'}$ soit
\printedpage{62}
% head: A. Grothendieck : Dérivateurs, Ch. XII, Édition des Universités de Montpellier II et Paris VII
% foot: Version 7 décembre 2002
dans $W$. Il faudra une notion plus forte, qui est suggérée par le cas où $T=\{t\}$ est la catégorie finale et contenue dans $\bar U$ (cas auquel nous avons essayé de nous ramener dans la proposition 7). Dans ce cas, on l’a dit (corollaire précédent), la condition que $i:U\longrightarrow\bar U=X$ soit dans $W$, équivaut à la $W$-asphéricité de $i$, i.e. à la $W$-asphéricité de $U/t$, et de même pour $i':U'\longrightarrow\overline{U'}=X'$ : elle est dans $W$ si et seulement si $U'/t'$ est $W$-asphérique, où
\msp{page 87}
$t'$ est l’unique élément de $T'(\simeq T)=X'\backslash U'$. Donc il faut des conditions sur $U/t$ qui assurent que non seulement $U/t$, mais aussi toutes les catégories similaires $U'/t$ \texttt{[plutôt} $U'/t'$\texttt{]}, soient $W$-asphériques. Dans le cas général où $T$ est quelconque, c’est cette condition-là que nous allons imposer aux $U/t$ ($t\in\Ob T$), et il s’avérera que cette fois-ci, elles sont stables \texttt{[plutôt elle est stable]} par tout cochangement de base $U\xrightarrow{p}U'$, i.e. que les $U'/t'$ ($t'\in\Ob T'$) satisferont alors à la même condition. C’est cette condition que nous allons à présent reformuler.

\subsection*{2.3 Objets totalement $W$-asphériques de \textnormal{Cat}, et morphismes totalement $W$-asphériques.}

Reprenons la situation
\[
U\overset{i}{\hookrightarrow}X\overset{j}{\hookleftarrow}T,\qquad\qquad\text{avec }T=\{t\}\text{, catégorie finale, }t\in\bar U\text{, i.e. }X=\bar U,
\]
et celle d’un cochangement de base $U\longrightarrow U'$, donnant naissance à un carré cocartésien dans Cat
\[
(*)\qquad\qquad
\begin{tikzcd}
U \arrow[r, hook, "i"] \arrow[d, "p"'] & X \arrow[d, "q"'] & T=\{t\} \arrow[l, hook', dotted, "j"'] \arrow[d, dotted, "\mathrm{iso}"] \\
U' \arrow[r, hook, "i'"] & X' & T'=\{t'\}. \arrow[l, hook', dotted, "j'"']
\end{tikzcd}
\]
La donnée de recollement $H$ entre $U$ et $T$, $H:\underbrace{U^\circ\times T}_{\simeq U^\circ}\longrightarrow\mathrm{Ens}$, peut s’interpréter comme un préfaisceau sur $U$, $H\in\Ob U^\wedge:U^\circ\longrightarrow\mathrm{Ens}$,
\msp{page 88}
et on a
\[
U/t\simeq U/H,
\]
où $U/H$ est interprété en termes du plongement canonique
\[
U\hookrightarrow U^\wedge.
\]
De même, le recollement entre $U'$ et $T'$ sera donné par un préfaisceau
\[
H'\in\Ob U^{\prime\wedge}:U^{\prime\circ}\longrightarrow\mathrm{Ens},
\]

\end{document}
